\documentclass[11pt,a4paper]{article}

\usepackage[utf8]{inputenc}
\usepackage[T1]{fontenc}
\usepackage{amsmath,amssymb}
\usepackage{graphicx}
\usepackage{booktabs}
\usepackage{longtable}
\usepackage{geometry}
\usepackage{hyperref}
\usepackage{enumitem}
\usepackage{xcolor}

\geometry{margin=2.7cm}
\hypersetup{colorlinks=true, linkcolor=blue!50!black, citecolor=blue!50!black, urlcolor=blue!50!black}

\newtheorem{hypothesis}{Hipótese}
\newtheorem{principle}{Princípio}

\title{\textbf{Prior-Guided Spectral Gating como Framework Geral de Atenção Espectral Interpretável:\\
Generalização entre Modalidades Espectroscópicas (NIR–Raman)}\\[0.3cm]
\large Proposta definitiva de projeto de pesquisa --- \texttt{pgsg\_2}\\
\normalsize Segundo projeto derivado do programa de pesquisa sobre Prior-Guided Spectral Gating (PGSG)}

\author{Aluno de Doutorado \and Clarimar José Coelho \\
\small Programa de Pós-Graduação, Pontifícia Universidade Católica de Goiás (PUC-Goiás)}

\date{\today}

\begin{document}

\maketitle

\begin{abstract}
\noindent
O projeto fundador do programa de pesquisa, \texttt{pgsg\_0}, propôs o \emph{Prior-Guided Spectral Gating} (PGSG), um mecanismo de atenção espectral guiado por priors quimiométricos, validado exclusivamente em espectroscopia no infravermelho próximo (NIR). O projeto \texttt{pgsg\_1} --- \emph{"Prior-Guided Spectral Gating with End-to-End Optimization for NIR Spectroscopy: Scalable Architecture for Interpretable Regression in Small-Sample Regimes"} --- atacou a primeira limitação declarada por \texttt{pgsg\_0} --- validação restrita a amostras pequenas --- demonstrando, com a base Mango DMC v3 ($n=1{,}448$ no protocolo intra-safra), que a arquitetura revisada PGSGv2 (gates sigmoid independentes, priors não-circulares, refinamento \emph{data-driven}) mantém vantagem preditiva sobre PLS mesmo em larga escala ($R^2=0{,}806$ vs.\ $0{,}568$). Este documento apresenta a proposta definitiva de \texttt{pgsg\_2}, que ataca a segunda limitação declarada: a validação restrita a uma única modalidade espectroscópica. Diferentemente de uma pré-proposta de brainstorming que cogitou generalizar simultaneamente para Raman, infravermelho médio (MIR) e espectrometria de massas (MS), esta proposta delimita o escopo empírico a duas modalidades --- NIR (reuso de \texttt{pgsg\_1}) e Raman (via o benchmark público RamanBench/\texttt{raman-data}) --- e formaliza um novo constructo, a \emph{topologia do gate espectral}, caracterizada por três descritores quantitativos: entropia, suavidade e esparsidade (índice de Hoyer). MIR e MS são explicitamente descartados do escopo empírico atual por razões de viabilidade e de adequação física do mecanismo, permanecendo como trabalho futuro do programa.
\end{abstract}

\tableofcontents

\section{Contexto no programa de pesquisa}

O projeto fundador \texttt{pgsg\_0} declara cinco limitações (Seção 6), cada uma endereçada por um projeto derivado:

\begin{center}
\begin{tabular}{cl}
\toprule
Projeto & Limitação atacada \\
\midrule
\texttt{pgsg\_0} & --- (artigo/projeto fundador; propõe o mecanismo PGSG) \\
\texttt{pgsg\_1} & Escalabilidade (validação restrita a $n=60$--$215$) \\
\textbf{\texttt{pgsg\_2}} & \textbf{Modalidade única (validação restrita a NIR)} \\
\texttt{pgsg\_3} & Alvo único (regressão multi-alvo) \\
\texttt{pgsg\_4} & Otimização sistemática de hiperparâmetros \\
\texttt{pgsg\_5} & Ausência de quantificação de incerteza \\
\bottomrule
\end{tabular}
\end{center}

\texttt{pgsg\_1} encerrou-se em estado de submissão (artigo \emph{"Prior-Guided Spectral Gating with End-to-End Optimization for NIR Spectroscopy: Scalable Architecture for Interpretable Regression in Small-Sample Regimes"}, formato \texttt{elsarticle}, periódico-alvo \emph{Chemometrics and Intelligent Laboratory Systems} --- Chemolab). \texttt{pgsg\_2} reaproveita a arquitetura PGSGv2 e a filosofia de contratos por estágio (operadores $\mathcal{G}, \mathcal{T}, \mathcal{P}, \mathcal{C}, \mathcal{M}, \mathcal{I}_{\text{pred}}, \mathcal{I}_{\text{interp}}, \Sigma$) já implementada e testada em \texttt{pgsg\_1}.

\section{Da pré-proposta à proposta definitiva}

A pré-proposta que originou este documento continha um insight central correto: reposicionar o PGSG como um \emph{operador geral de modulação espectral guiado por conhecimento físico/químico}, em vez de uma técnica específica para NIR. Este reposicionamento é mantido. Três decisões de escopo, porém, foram revisadas:

\begin{enumerate}[label=(\alph*)]
\item \textbf{Redução de quatro para duas modalidades.} A pré-proposta cogitava NIR, Raman, MIR e MS na mesma investigação. Isso reproduziria, em sentido inverso, o problema que \texttt{pgsg\_1} corrigiu: uma validação cujo escopo declarado excede o que pode ser executado e testado com rigor dentro do prazo de um projeto. Raman é escolhida como segunda modalidade por (i) possuir um benchmark público de grande escala e recém-disponibilizado (RamanBench, Seção~\ref{sec:dados}), e (ii) apresentar uma estrutura física de pico (bandas estreitas, alta razão sinal-ruído localizada) que contrasta de forma nítida com as bandas largas e sobrepostas do NIR --- o contraste mais informativo para testar a hipótese de topologia de gate.
\item \textbf{Exclusão de espectrometria de massas do escopo empírico.} O argumento físico que justifica o gating --- correlação entre posições espectrais adjacentes ao longo de um eixo contínuo --- não se sustenta em MS, cujos sinais são discretos em razão massa/carga sem a mesma estrutura de vizinhança local. Incluir MS sem adaptar a arquitetura seria uma extensão fraca. MS permanece como direção de trabalho futuro, condicionada a uma reformulação do operador de gating para dados não contínuos.
\item \textbf{Formalização matemática dos descritores de topologia do gate.} As expressões de entropia, suavidade e esparsidade na pré-proposta precisam de correção: gates sigmoid independentes não somam 1 (pré-requisito para entropia de Shannon), e a razão de esparsidade proposta não corresponde a nenhuma medida padrão da literatura de processamento de sinais. A Seção~\ref{sec:topologia} apresenta as definições corrigidas.
\end{enumerate}

\section{Pergunta científica e hipóteses}

\textbf{Pergunta central:} O mecanismo PGSGv2, validado em NIR, generaliza --- sem alteração arquitetural --- para uma segunda modalidade espectroscópica (Raman), e a topologia do gate aprendido difere sistematicamente entre as duas modalidades de forma consistente com a física subjacente de cada uma?

\begin{hypothesis}[Generalização sem mudança arquitetural]
PGSGv2, aplicado a dados Raman sem modificação da arquitetura (apenas reconfiguração do prior de inicialização), supera PLS em pelo menos um cenário de regressão do RamanBench, replicando qualitativamente o resultado obtido em NIR ($\Delta R^2_{\text{PLS}} > 0$).
\end{hypothesis}

\begin{hypothesis}[Esparsidade diferencial]
O gate aprendido em dados Raman apresenta esparsidade (índice de Hoyer) significativamente maior do que o gate aprendido em dados NIR, refletindo a natureza de picos estreitos do espalhamento Raman frente às bandas largas e sobrepostas de sobretons/combinações no NIR.
\end{hypothesis}

\begin{hypothesis}[Suavidade diferencial]
O gate aprendido em NIR apresenta suavidade (variação total normalizada) significativamente maior do que em Raman, i.e., o gate varia mais abruptamente entre bandas vizinhas no caso Raman.
\end{hypothesis}

\begin{hypothesis}[Vantagem do prior fisicamente motivado]
A inicialização do gate por atribuições vibracionais da literatura (não derivadas dos dados de treino) produz desempenho preditivo igual ou superior, e maior estabilidade do gate entre sementes de treinamento, em comparação com inicialização não informada, em ambas as modalidades.
\end{hypothesis}

\begin{hypothesis}[Correlação estrutura--física]
A entropia e a suavidade do gate aprendido correlacionam-se com a largura de banda espectral conhecida das transições vibracionais relevantes para o analito-alvo em cada modalidade.
\end{hypothesis}

\section{Contribuições}

\begin{enumerate}
\item Extensão do mecanismo PGSGv2 para uma segunda modalidade espectroscópica pública (Raman), com reuso do pipeline de oito operadores de \texttt{pgsg\_1} e mudança mínima de arquitetura.
\item Formalização de três descritores quantitativos da estrutura do gate aprendido --- entropia, suavidade, esparsidade --- como caracterização reutilizável de \emph{topologia do gate espectral}.
\item Teste falsificável da hipótese de que a topologia do gate reflete a física da modalidade, e não apenas o conjunto de dados específico.
\item Avaliação da robustez do princípio anti-circularidade de prior (estabelecido em \texttt{pgsg\_1}/ADR correspondente) em um segundo domínio espectroscópico.
\item Extensão dos contratos de estágio (\emph{ingestion}, \emph{priors}) do pipeline para suportar múltiplas modalidades, preparando a base de código para os projetos subsequentes do programa.
\end{enumerate}

\section{Bases de dados}
\label{sec:dados}

\begin{center}
\begin{tabular}{p{2.6cm}p{5.6cm}p{6.7cm}}
\toprule
Modalidade & Base & Papel no projeto \\
\midrule
NIR & Mango DMC v3 (reuso de \texttt{pgsg\_1}); Tecator (âncora de domínio) & Ramo de referência já validado; nenhum reprocessamento de arquitetura necessário. \\
Raman & \texttt{bioprocess\_substrates} (Lange et al., 2026, \emph{Measurement}), acessado via RamanBench/\texttt{raman-data} --- $n=6.472$ espectros reais válidos (6.960 originais menos 488 sem rótulo de glicose), 1.870 bandas, eixo 390,63--3384,70 cm$^{-1}$; alvo: concentração de glicose & Fonte primária do ramo Raman, definida em ADR-0001 (\texttt{pgsg\_2}) e confirmada por execução real na máquina com GPU (16/07/2026). \\
\bottomrule
\end{tabular}
\end{center}

\noindent\textbf{Nota sobre escala.} Diferentemente da pré-proposta original --- que temia um regime de amostra tipicamente pequeno em Raman (mediana de 235 espectros entre os 58 datasets de regressão do RamanBench, 87\% abaixo de 500 amostras) --- o dataset selecionado (\texttt{bioprocess\_substrates}, $n=6{.}960$) está na mesma ordem de grandeza do ramo NIR (Mango DMC v3, $n\approx 12{.}000$), o que permite testar as Hipóteses 1--5 sem o fator de confusão de escalas de amostra muito diferentes entre modalidades.

\noindent\textbf{Nota sobre execução.} O carregamento real (fonte Hugging Face: \texttt{chlange/SubstrateMixRaman}) foi executado e validado na máquina com GPU/60~GB RAM em 16/07/2026, confirmando $n=6.472$ amostras válidas, 1.870 bandas e a presença inequívoca do alvo "glucose" em \texttt{target\_names} --- ver ADR-0001 (\texttt{pgsg\_2}).

\section{Pré-processamento}

\begin{center}
\begin{tabular}{p{2.6cm}p{11.5cm}}
\toprule
Modalidade & Etapas de pré-processamento ($\mathcal{T}$) \\
\midrule
NIR & Reuso integral do operador $\mathcal{T}$ de \texttt{pgsg\_1}: remoção de bandas zeradas por safra (ADR-0005), SNV/derivada conforme protocolo já validado. \\
Raman & (i) Correção de linha de base (ex.: \emph{asymmetric least squares}, ALS), necessária pela fluorescência de fundo característica de Raman; (ii) remoção de \emph{cosmic rays} (picos espúrios de curta largura); (iii) normalização de intensidade (vetor unitário ou SNV); (iv) alinhamento de eixo de deslocamento Raman (cm$^{-1}$) entre instrumentos, quando aplicável ao dataset selecionado. \\
\bottomrule
\end{tabular}
\end{center}

Os contratos do operador $\mathcal{T}$ precisam ser estendidos (não reescritos) para aceitar um parâmetro de modalidade que seleciona a cadeia de pré-processamento apropriada, mantendo a mesma interface (\texttt{fit\_transform} retornando \texttt{(X\_ndarray, y\_ndarray)}) já usada em \texttt{pgsg\_1}.

\section{Métodos}

\subsection{Arquitetura}
Reuso integral da PGSGv2 (gates sigmoid independentes; regularização KL entre gate aprendido e prior; refinamento \emph{data-driven}). Nenhuma mudança arquitetural é introduzida entre modalidades --- esta é, precisamente, a condição de teste da Hipótese~1.

\subsection{Priors por modalidade}
\begin{itemize}
\item \textbf{NIR:} prior VIP literatura-informado, conforme já estabelecido em \texttt{pgsg\_1} (evitando a inicialização circular diagnosticada naquele projeto).
\item \textbf{Raman:} prior construído a partir de atribuições vibracionais conhecidas da literatura para o(s) analito(s)-alvo do dataset selecionado (ex.: estiramento C=O carboxílico $\approx 1700$~cm$^{-1}$ para ácidos orgânicos), nunca derivado dos dados de treino do próprio experimento.
\end{itemize}

\subsection{Protocolo experimental}
Protocolo intra-domínio (\emph{within-dataset}), replicando a decisão de \texttt{pgsg\_1} de evitar comparações entre domínios heterogêneos sem validação prévia (cf.\ degradação observada no protocolo cross-safra). Conjunto de teste fixo, disjunto e de tamanho constante por dataset.

\subsection{Topologia do gate espectral}
\label{sec:topologia}

Seja $g \in (0,1)^p$ o vetor de gate aprendido sobre $p$ bandas/números de onda.

\begin{align}
\text{Entropia normalizada:}\quad & H(g) = -\sum_{i=1}^{p} \hat{g}_i \log \hat{g}_i, \qquad \hat{g}_i = \frac{g_i}{\sum_{j=1}^p g_j} \\
\text{Suavidade (variação total normalizada):}\quad & \mathrm{TV}(g) = \frac{1}{p-1}\sum_{i=1}^{p-1} |g_i - g_{i+1}| \\
\text{Esparsidade (índice de Hoyer):}\quad & S(g) = \frac{\sqrt{p} - \|g\|_1/\|g\|_2}{\sqrt{p} - 1} \in [0,1]
\end{align}

$S(g)=0$ corresponde ao gate uniformemente denso; $S(g)=1$ corresponde a um gate \emph{one-hot}. Esta formalização substitui a razão $\|g\|_2/\|g\|_1$ da pré-proposta original, que não é uma medida de esparsidade padrão nem tem direção de interpretação definida.

\subsection{Baselines adicionais para H1}
\label{sec:baselines}

A comparação central de H1 permanece PGSGv2 vs.\ PLS. Para aproximar a seção de resultados do estado da arte internacional em atenção espectral, o conjunto de baselines é ampliado com três referências suplementares de desempenho preditivo (não de topologia de gate):

\begin{itemize}
\item \textbf{ABRN} (\emph{Attention-based Residual Network}, Huang et al., 2019) --- rede residual com módulo de atenção aplicada diretamente a espectros NIR brutos, sem seleção manual de variáveis; isola o efeito do prior quimiométrico do PGSGv2 frente a atenção não guiada.
\item \textbf{SpecFuseNet} (Kabiso \& Ko, 2025) --- autoencoder residual com atenção de canal (FusedECA) e gate residual espectral (SRG), originalmente proposto para classificação de cultivares por NIR, adaptado aqui para regressão.
\item \textbf{Baselines lineares esparsos} (Elastic Net, Sparse PLS) --- alternativas de seleção de variáveis com interpretabilidade nativa, citadas na literatura de interpretabilidade em NIR; funcionam como ponte entre o baseline linear puro (PLS) e os baselines de atenção profunda.
\end{itemize}

Esses baselines não alteram o escopo de H1--H5 --- continuam centradas em PGSGv2 vs.\ PLS --- e entram apenas como comparação suplementar de desempenho no artigo resultante.

\subsection{Nota: trilha de imageamento hiperespectral (HSI)}
\label{sec:hsi-nota}

Uma proposta de financiamento em preparação (CNPq/FNDCT, Chamada Universal 06/2026, Faixa C) associada a este projeto prevê a aquisição de infraestrutura própria --- câmera HSI-SWIR e servidor com GPU dedicado --- para reduzir a dependência de acesso pontual a equipamento de terceiros e abrir uma terceira frente de validação (dados espacialmente resolvidos). Essa frente é tratada explicitamente como \textbf{piloto de infraestrutura, fora do escopo empírico das Hipóteses H1--H5} deste documento: um cubo HSI pode ser tratado como uma coleção de espectros pontuais (um por pixel), permitindo reaproveitar os operadores $\mathcal{G}/\mathcal{T}/\mathcal{P}/\mathcal{M}$ já validados aqui, com extensão natural do operador $\mathcal{I}_{\text{interp}}$ para mapas espaciais de gate --- direção que alimenta \texttt{pgsg\_3} (multi-alvo), não uma hipótese testada em \texttt{pgsg\_2}.

\section{Riscos e limitações}

\begin{itemize}
\item \textbf{Escopo de modalidade:} MIR e MS não fazem parte do escopo empírico deste projeto (Seção~2); resultados não devem ser generalizados além de NIR/Raman sem validação adicional.
\item \textbf{Tamanho amostral em Raman:} datasets do RamanBench majoritariamente pequenos; comparações estatísticas entre modalidades devem reportar intervalos de confiança e não apenas estimativas pontuais de $R^2$.
\item \textbf{Comparabilidade entre eixos espectrais:} NIR (nm) e Raman (cm$^{-1}$) não compartilham unidade física direta; os descritores de topologia (Seção~\ref{sec:topologia}) são adimensionais por construção para permitir comparação, mas a interpretação física de "banda vizinha" difere entre as duas escalas e deve ser discutida explicitamente na seção de resultados.
\item \textbf{Dataset Raman definido (ADR-0001).} A seleção foi feita por inspeção de metadados, sem download efetivo; o carregamento real do \texttt{bioprocess\_substrates} e a confirmação de que $n=6{.}960$ se mantém após limpeza de dados ficam como primeira etapa de execução na máquina com GPU.
\end{itemize}

\section{Cronograma}

\begin{center}
\begin{tabular}{p{1.6cm}p{11.5cm}}
\toprule
Meses & Atividade \\
\midrule
1--2 & Escopo de dados: inspeção do \texttt{raman-data} e seleção do dataset Raman \textbf{(concluído --- ADR-0001, \texttt{bioprocess\_substrates})}; extensão do contrato $\mathcal{G}$ (ingestão) para múltiplas modalidades (em andamento). \\
3--4 & Extensão do operador $\mathcal{T}$ (pré-processamento específico de Raman: ALS, remoção de \emph{cosmic rays}); testes de contrato. \\
5--6 & Definição do prior fisicamente motivado para Raman (operador $\mathcal{P}$); execução do protocolo experimental intra-domínio para NIR e Raman com PGSGv2 sem alteração arquitetural (operadores $\mathcal{C}, \mathcal{M}, \mathcal{I}_{\text{pred}}$). \\
7--8 & Implementação e cálculo dos descritores de topologia do gate (entropia, suavidade, esparsidade) no operador $\mathcal{I}_{\text{interp}}$; testes estatísticos das Hipóteses 2--5. \\
9--10 & Síntese de resultados ($\Sigma$): figuras, tabelas, testes de hipótese consolidados; execução dos baselines suplementares (ABRN, SpecFuseNet, Elastic Net, Sparse PLS --- Seção~\ref{sec:baselines}) para a comparação de desempenho preditivo. \\
11 & Redação do artigo (formato \texttt{elsarticle}, Chemolab), revisão interna. \\
12 & Submissão e resposta a eventuais revisões internas antes de envio. \\
\bottomrule
\end{tabular}
\end{center}

\appendix

\section{Pipeline proposto (detalhamento)}

O pipeline mantém a composição estabelecida em \texttt{pgsg\_1}:
\[
\Pi : \mathcal{D} \times \Theta \longrightarrow \mathcal{R}, \qquad
\Pi = \Sigma \circ \mathcal{I}_{\text{interp}} \circ \mathcal{I}_{\text{pred}} \circ \mathcal{M} \circ \mathcal{C} \circ \mathcal{P} \circ \mathcal{T} \circ \mathcal{G}
\]

\begin{center}
\begin{longtable}{p{1.3cm}p{2.6cm}p{9.8cm}}
\toprule
Operador & Submódulo & Extensão necessária para \texttt{pgsg\_2} \\
\midrule
\endhead
$\mathcal{G}$ & \texttt{ingestion} & Adicionar carregador para dataset(s) Raman via \texttt{raman-data}; contrato \texttt{SpectralDataset} ganha valor de \texttt{metadata[domain]} = \texttt{"raman"}, preservando \texttt{wavelength\_unit} = cm$^{-1}$. \\
$\mathcal{T}$ & \texttt{preprocessing} & Nova cadeia de pré-processamento Raman (ALS, remoção de \emph{cosmic rays}, normalização); interface \texttt{fit\_transform} inalterada. \\
$\mathcal{P}$ & \texttt{priors} & Novo tipo de prior: atribuição vibracional literatura-informada (não-VIP), com mesma assinatura de \texttt{VIPPrior.fit(X, y)} por compatibilidade de contrato. \\
$\mathcal{C}$ & \texttt{scenarios} & Reuso direto de \texttt{ScenarioGenerator}; protocolo intra-domínio (sem cenário cross-modalidade neste projeto). \\
$\mathcal{M}$ & \texttt{models} & Reuso integral de PGSGv2, PLS; nenhuma mudança arquitetural entre modalidades (condição da Hipótese 1). \\
$\mathcal{I}_{\text{pred}}$ & \texttt{evaluation} & Reuso integral; $\Delta R^2_{\text{PLS}}$ calculado por modalidade. \\
$\mathcal{I}_{\text{interp}}$ & \texttt{interpretability} & Nova funcionalidade: cálculo de $H(g)$, $\mathrm{TV}(g)$, $S(g)$ (Seção~\ref{sec:topologia}) por modalidade e por semente de treinamento. \\
$\Sigma$ & \texttt{synthesis} & Testes de hipótese (Hipóteses 2--5) comparando distribuições de $H(g)$, $\mathrm{TV}(g)$, $S(g)$ entre NIR e Raman. \\
\bottomrule
\end{longtable}
\end{center}

\section{Referências indicativas}

\begin{itemize}
\item Koddenbrock, M., Lange, C., Legner, R., Jaeger, M., Kögler, M., et al. \emph{RamanBench: A Large-Scale Benchmark for Machine Learning on Raman Spectroscopy}. arXiv:2605.02003, 2026.
\item Echtermeyer, A., Marks, C., Mitsos, A., Viell, J. \emph{Inline Raman spectroscopy and indirect hard modeling for concentration monitoring of dissociated acid species}. Applied Spectroscopy, 75(5), 506--519, 2021.
\item Béjar-Grimalt, J., Sánchez-Illana, Á., Quintás, G., Byrne, H.~J., Pérez-Guaita, D. \emph{Monte Carlo peaks: simulated datasets to benchmark machine learning algorithms for clinical spectroscopy}. Chemometrics and Intelligent Laboratory Systems, 2025.
\item Anderson, N.~T., Walsh, K.~B., Subedi, P.~P. \emph{Mango DMC v3}. Mendeley Data, DOI: 10.17632/46htwnp833.3.
\item Huang, L., Guo, S., Wang, Y., Wang, S., Chu, Q., Li, L., Bai, T. \emph{Attention based residual network for medicinal fungi near infrared spectroscopy analysis}. Mathematical Biosciences and Engineering, 16(4), 3003--3017, 2019.
\item Kabiso, A.~C., Ko, C.-H. \emph{Attention-enhanced residual autoencoder for NIR spectral feature extraction and classification of grain varieties}. Scientific Reports, 15, 32750, 2025.
\end{itemize}

\end{document}
